% !TEX program = xelatex
\documentclass[11pt]{ctexart}
\usepackage[a4paper,margin=2.2cm,top=2cm,bottom=2cm]{geometry}
\usepackage{amsmath,amssymb,amsthm,bm}
\usepackage[table]{xcolor}
\usepackage{enumitem,array,booktabs,lastpage,fancyhdr}
\usepackage{tcolorbox}
\tcbuselibrary{skins,breakable}
\usepackage{hyperref}
\hypersetup{hidelinks,pdftitle={线性代数期中模拟卷 2021}}

\definecolor{acc}{HTML}{1B4F8A}

\xeCJKsetup{CJKmath=true}
\setlength{\parindent}{0pt}
\setlength{\parskip}{6pt}
\setlist[enumerate]{leftmargin=2em,itemsep=10pt,topsep=6pt}

\newcommand{\Rr}{\mathbb{R}}
\newcommand{\rank}{\operatorname{rank}}
\newcommand{\Nul}{\operatorname{Nul}}
\newcommand{\TT}{^{\mathsf T}}
\newcommand{\bv}[1]{\begin{pmatrix}#1\end{pmatrix}}

\pagestyle{fancy}\fancyhf{}
\fancyhead[L]{\small\color{acc}复旦大学信息科学与工程学院 \quad《线性代数》期中考试}
\fancyhead[R]{\small\color{acc}第 \thepage\ 页 \,/\, 共 \pageref*{LastPage} 页}
\renewcommand{\headrulewidth}{0.4pt}
\fancyfoot[C]{\small 2021 年 11 月 \quad 闭卷 \quad 100 分钟}

\begin{document}

\begin{center}
{\Large\bfseries 复旦大学信息科学与工程学院}\\[4pt]
{\Large\bfseries 《线性代数》期中考试试卷}\\[8pt]
{\small 课程代码 COMP120004.08 \quad 闭卷 \quad 2021 年 11 月 \quad 满分 100 分 \quad 答题时间 100 分钟}
\end{center}

\vspace{0.2cm}
\begin{center}
\begin{tabular}{|c|c|c|c|c|c|c|c|c|c|}
\hline
题号 & 一 & 二 & 三 & 四 & 五 & 六 & 七 & 八 & 九 \\
\hline
分值 & 10 & 12 & 16 & 10 & 10 & 10 & 10 & 10 & 12 \\
\hline
得分 & \phantom{00} & \phantom{00} & \phantom{00} & \phantom{00} & \phantom{00} & \phantom{00} & \phantom{00} & \phantom{00} & \phantom{00} \\
\hline
\end{tabular}
\end{center}

\vspace{0.2cm}
{\small 专业 \underline{\hspace{3cm}} \quad 学号 \underline{\hspace{3cm}} \quad 姓名 \underline{\hspace{3cm}} \quad 成绩 \underline{\hspace{2cm}}}

\vspace{0.3cm}
{\small\textbf{注意}：答案必须写在试卷上，做在草稿纸上无效。}

\vspace{0.4cm}
\hrule
\vspace{0.4cm}

\begin{enumerate}

%-----------------------------------------
\item （10 分）若 $\mathbf x=[1,-2,1,0]\TT$，$\mathbf y=[0,1,-2,1]\TT$，求 $\mathbf x\TT\mathbf y$、$\mathbf x\mathbf y\TT$ 与 $(\mathbf x\mathbf y\TT)^3$ 的结果，要求有计算过程。

\vspace{7cm}

%-----------------------------------------
\item （12 分）讨论以下线性方程组当 $\lambda$ 取不同值时的解的情况（无解、唯一解、无穷多解），若有解请写出解集。
\[
\begin{cases}
\lambda x_1+x_2+x_3=1\\
x_1-\lambda x_2+x_3=\lambda\\
x_1+x_2-\lambda x_3=\lambda^2
\end{cases}
\]

\newpage

%-----------------------------------------
\item （16 分）已知在常规矩阵加法和数乘规则下，所有 $2\times2$ 的矩阵构成的集合 $M^{2\times2}$ 为线性空间。

\begin{enumerate}[label=(\arabic*),leftmargin=2.2em,itemsep=8pt]
\item 请证明所有 $2\times2$ 的对称矩阵（即转置不变的矩阵）构成的空间 $S^{2\times2}$ 为线性子空间；（4 分）
\item 给出 $S^{2\times2}$ 的两组不同的基；（4 分）
\item 计算这两组基之间的过渡矩阵；（4 分）
\item 计算矩阵 $\bv{2&1\\1&2}$ 在这两组基下的展开系数。（4 分）
\end{enumerate}

\vspace{8cm}

%-----------------------------------------
\item （10 分）请计算三维空间中坐标为 $(1,-1,1)$ 的点 $\mathbf x$ 到矩阵
$A=\bv{1&0&1\\0&1&2\\0&0&0}$ 的零空间 $\Nul(A)$ 的距离，要求有具体步骤。

\newpage

%-----------------------------------------
\item （10 分）设 $A=\bv{1&0\\0&0}$，$B=\bv{0&1\\0&0}$，求两个零空间之和 $\Nul(A)+\Nul(B)$ 的基与维数。

\vspace{6cm}

%-----------------------------------------
\item （10 分）设 $\mathbf v_1,\mathbf v_2,\mathbf v_3,\mathbf v_4\in\Rr^4$ 为四个线性无关向量，
请证明 $\dim\operatorname{span}\{\mathbf v_1-\mathbf v_2,\ \mathbf v_2-\mathbf v_3,\ \mathbf v_3-\mathbf v_4,\ \mathbf v_4-\mathbf v_1\}=3$。

\newpage

%-----------------------------------------
\item （10 分）对于矩阵 $A_{m\times n}$，$B_{n\times s}$，$C_{s\times t}$，请证明 Frobenius 不等式
\[
\rank(ABC)\ge\rank(AB)+\rank(BC)-\rank(B),
\]
并说明它与 Sylvester 不等式 $\rank(A)+\rank(B)-n\le\rank(AB)$ 的关系。

\vspace{7cm}

%-----------------------------------------
\item （10 分）如果 $W$ 是线性空间 $V$ 的线性子空间，请证明：

\begin{enumerate}[label=(\arabic*),leftmargin=2.2em,itemsep=8pt]
\item $W$ 的零元即为 $V$ 中的零元；（2 分）
\item 若 $\dim W=\dim V$，则 $W=V$。（8 分）
\end{enumerate}

\newpage

%-----------------------------------------
\item （12 分）$W$ 是线性空间 $V$ 的线性子空间，已知 $\dim V=n>\dim W=r$。
若给定 $W$ 的一组基 $\{\mathbf w_1,\mathbf w_2,\dots,\mathbf w_r\}$，
请证明：可以找到 $n-r$ 个线性无关向量 $\{\mathbf v_1,\mathbf v_2,\dots,\mathbf v_{n-r}\}$，
使得它们与 $\{\mathbf w_1,\mathbf w_2,\dots,\mathbf w_r\}$ 共同构成 $V$ 的基。

\vspace{10cm}

\end{enumerate}

\vfill
\begin{center}
{\small —— 试卷结束 ——}
\end{center}

\end{document}
